\chapter{Logica dei predicati}\label{cap:predicati}

La logica proposizionale non basta per fare matematica: non sa parlare di \textbf{oggetti} (numeri, insiemi, \dots)
né dire «\textbf{per ogni}» o «\textbf{esiste}». La \textbf{logica dei predicati} (o \textbf{logica del primo ordine})
aggiunge proprio questi ingredienti. Seguiremo lo stesso percorso del capitolo precedente: alfabeto, sintassi
(termini e formule ben formate) e infine semantica.

% ──────────────────────────────────────────────────────────────
\section{Il linguaggio dell'aritmetica}\label{sec:aritmetica}

Informalmente, il modello che abbiamo per obiettivo è l'\textbf{aritmetica}, cioè il linguaggio del primo ordine
\[ (\mbN, +, \cdot, 0, s, \leq, =) \]
dove:
\begin{itemize}
    \item $\mbN$ è l'\textbf{insieme} dei numeri naturali, cioè gli «oggetti» di cui vogliamo parlare;
    \item $+$ e $\cdot$ sono \textbf{operazioni binarie};
    \item $0$ è una \textbf{costante};
    \item $s$ è una \textbf{funzione unaria}, il \textit{successore}: $s \colon n \in \mbN \mapsto n + 1 \in \mbN$;
    \item $\leq$ è un \textbf{predicato} (binario);
    \item $=$ è l'\textbf{uguale}.
\end{itemize}
Vogliamo poter scrivere cose come $x \leq y$ oppure $(x = y) \rightarrow \big(\forall z\,(z + x = 0)\big)$.

\dfn[arieta]{Arietà}{
    L'\textbf{arietà} (o $i$-arietà) di una funzione o di un predicato è il \textbf{numero dei suoi argomenti}.
    Una funzione o un predicato con $n$ argomenti si dice \textbf{$n$-ario} (\textit{unario} se $n = 1$,
    \textit{binario} se $n = 2$).
}

Nell'aritmetica, $+$ è una funzione binaria (o «2-aria», di arietà 2), la funzione $s$ è una funzione unaria,
mentre $\leq$ è un predicato binario, e così via.

% ──────────────────────────────────────────────────────────────
\section{L'alfabeto}\label{sec:alfabeto-pred}

\dfn[alfabeto-pred]{Alfabeto della logica dei predicati}{
    L'\textbf{alfabeto} della logica dei predicati è
    \[ \mcA = (V,\ \mathrm{Cost},\ P,\ F,\ B,\ \mathrm{Conn},\ Q,\ \{=\}) \]
    dove:
    \begin{itemize}
        \item $V = \{x_1, \dots, x_n\}$ è l'insieme delle \textbf{variabili};
        \item $\mathrm{Cost} = \{c_1, \dots, c_n\}$ è l'insieme delle \textbf{costanti} (individuali);
        \item $P = \{p_1^{i_1}, \dots, p_n^{i_n}\}$ è l'insieme delle \textbf{lettere predicative} (o \textbf{relazionali}) con $i$-arietà;
        \item $F = \{f_1^{i_1}, \dots, f_n^{i_n}\}$ è l'insieme delle \textbf{lettere funzionali} con $i$-arietà;
        \item $B = \{(,\ )\} \cup \{,\}$ è l'insieme degli \textbf{accessori} (parentesi e virgola);
        \item $\mathrm{Conn} = \{\neg, \wedge, \vee, \rightarrow, \leftrightarrow\}$ è l'insieme dei \textbf{connettivi logici};
        \item $Q = \{\forall\}$ è l'insieme dei \textbf{quantificatori};
        \item $\{=\}$ contiene il simbolo di \textbf{uguaglianza}.
    \end{itemize}
}

L'apice di una lettera ne indica l'arietà: ad esempio $p_1^{2}$ è una lettera predicativa binaria e $f_3^{1}$ è
una lettera funzionale unaria. Nell'aritmetica: $0$ è una costante, $s$, $+$ e $\cdot$ sono lettere funzionali
(di arietà 1, 2 e 2) e $\leq$ è una lettera predicativa binaria.

\warning{Attenzione ai nomi}{
    Qui $F$ è l'insieme delle \textbf{lettere funzionali}: non va confuso con il valore di verità $\falso$
    né con l'insieme $\mcF$ delle FBF del Capitolo~1. Allo stesso modo, $P$ qui è l'insieme delle lettere predicative
    e non quello delle variabili proposizionali.
}

% ──────────────────────────────────────────────────────────────
\section{Sintassi: i termini}\label{sec:termini}

Sintassi, ovvero: quali formule ha senso scrivere? Partiamo dai \textbf{termini}, che rappresentano gli
«oggetti» di cui parliamo (nell'aritmetica: i numeri).

\dfn[termini]{Termini}{
    I \textbf{termini} sono definiti dalle seguenti regole:
    \begin{enumerate}
        \item ogni variabile è un termine;
        \item ogni costante è un termine;
        \item se $f$ è una lettera funzionale $n$-aria e $t_1, \dots, t_n$ sono termini, allora $f(t_1, \dots, t_n)$ è un termine.
    \end{enumerate}
}

\begin{esempio}{Termini dell'aritmetica}
    Sono termini: $0$, \ $x$, \ $s(0)$, \ $s(s(0))$, \ $+(x, s(0))$, che scriveremo più comodamente $x + s(0)$.

    Non sono termini: $s(0, x)$ (perché $s$ è unaria e le si danno due argomenti) e $x \leq 0$
    (perché $\leq$ è una lettera \textbf{predicativa}: $x \leq 0$ non è un oggetto, ma un'affermazione, cioè una formula).
\end{esempio}

% ──────────────────────────────────────────────────────────────
\section{Sintassi: le formule ben formate}\label{sec:fbf-pred}

\dfn[fbf-pred]{Formule ben formate (FBF) del primo ordine}{
    Le \textbf{formule ben formate} della logica dei predicati sono definite dalle seguenti regole:
    \begin{enumerate}
        \item se $p$ è una lettera predicativa $n$-aria e $t_1, \dots, t_n$ sono termini, allora $p(t_1, \dots, t_n)$
              è una FBF, detta \textbf{atomica} (ad esempio $0 \leq s(0)$);
        \item se $t_1$ e $t_2$ sono termini, allora $t_1 = t_2$ è una FBF;
        \item se $a$ e $b$ sono FBF, allora $\neg a$, \ $a \wedge b$, \ $a \vee b$, \ $a \rightarrow b$, \ $a \leftrightarrow b$ sono FBF;
        \item se $a$ è una FBF e $x$ è una variabile, allora $(\forall x)(a)$ è una FBF.
              La FBF $a$ si dice \textbf{scopo} del quantificatore $\forall x$.
    \end{enumerate}
}

\info{Lettere predicative come proprietà}{
    Se $p$ è una lettera predicativa e $x$ una variabile, scriveremo $p(x)$ stando a significare che
    «\textbf{$x$ ha la proprietà $p$}».
}

\textbf{Significato di $\forall$.} La FBF $(\forall x)(a)$ vuol dire che $a$ (vista come un predicato su $x$) vale
\textbf{per ogni} termine $t$ messo al posto di $x$. Ad esempio $(\forall x)(0 \leq x)$ dice che ogni numero naturale
è maggiore o uguale a $0$.

Notiamo che $a$ può anche \textbf{non contenere} $x$: in quel caso $(\forall x)(a)$ equivale ad $a$.
Esempio: $\forall x\big((s(0) < s(s(0))) \wedge (2 = s(1))\big)$ dice semplicemente che $s(0) < s(s(0))$ e $2 = s(1)$.

\info{Abbreviazioni usate negli esempi}{
    Per comodità scriveremo:
    \begin{itemize}
        \item $t_1 \neq t_2$ al posto di $\neg(t_1 = t_2)$;
        \item $t_1 < t_2$ al posto di $(t_1 \leq t_2) \wedge (t_1 \neq t_2)$, e analogamente $\geq$ e $>$;
        \item $1$ al posto di $s(0)$, $2$ al posto di $s(s(0))$, e così via;
        \item $xy$ al posto di $x \cdot y$.
    \end{itemize}
}

% ──────────────────────────────────────────────────────────────
\section{Quantificatore esistenziale ed esistenziale unico}\label{sec:esiste}

L'unico quantificatore dell'alfabeto è $\forall$. A questa lista aggiungiamo altri due quantificatori,
\textbf{definiti} a partire da quello che abbiamo.

\dfn[esiste]{Quantificatore esistenziale}{
    Sia $a$ una FBF e $x$ una variabile. Definiamo
    \[ (\exists x)(a) \ \defsse\ \neg\big((\forall x)(\neg a)\big). \]
    Si legge «\textbf{esiste} $x$ tale che $a$»: non è vero che, per ogni $x$, vale $\neg a$.
}

\dfn[esiste-unico]{Quantificatore esistenziale unico}{
    Sia $p(x)$ una FBF. Definiamo
    \[ (\exists! x)\big(p(x)\big) \ \defsse\ (\exists x)\Big((\forall y)\big(p(y) \leftrightarrow x = y\big)\Big) \]
    dove $x$ e $y$ sono due variabili tali che $x \neq y$ e $y$ non compare in $p(x)$.
    Si legge «\textbf{esiste un unico} $x$ tale che $p(x)$».
}

Perché la definizione di $\exists!$ funziona? Dice che esiste un $x$ tale che, per ogni $y$, $p(y)$ vale
\textbf{esattamente quando} $y = x$. Prendendo $y = x$ si ottiene che $x$ ha la proprietà $p$; e se un qualunque
$y$ ha la proprietà $p$, allora $y = x$. Cioè: $x$ ha la proprietà $p$ ed è l'unico ad averla.

\info{Notazione}{
    Invece di $(\forall x)(f)$ scriveremo anche $\forall x(f)$, e così anche per $\exists$ ed $\exists!$.
}

% ──────────────────────────────────────────────────────────────
\section{Variabili libere e vincolate}\label{sec:libere}

\dfn[libere-vincolate]{Variabili libere e vincolate}{
    Sia $f$ una FBF. Un'occorrenza di una variabile $x$ si dice \textbf{vincolata} in $f$ se appare nella forma
    «$\forall x$» (e quindi anche «$\exists x$») oppure se si trova nello \textbf{scopo} di un quantificatore
    $\forall x$ (o $\exists x$) all'interno di $f$.

    Un'occorrenza che non è nello scopo di nessun quantificatore sulla stessa variabile si dice \textbf{libera}.
}

\begin{esempio}{Occorrenze libere e vincolate}
    \begin{itemize}
        \item In $\forall x\,(x < 0)$ sia la prima sia la seconda occorrenza di $x$ sono vincolate.
        \item In $\exists x\,\big(p(x) \rightarrow p(y)\big)$, $x$ è sempre vincolata, mentre $y$ è libera.
    \end{itemize}
\end{esempio}

\dfn[chiusa]{Formula chiusa}{
    Una FBF si dice \textbf{chiusa} quando non ha variabili libere (cioè è priva di occorrenze libere di variabili).
}

% ──────────────────────────────────────────────────────────────
\section{Cenni di semantica: formule valide}\label{sec:valide}

Per quanto riguarda la semantica della logica del primo ordine, diversamente da quanto accadeva nella logica
proposizionale, ci troviamo davanti a FBF che possono \textbf{non avere alcun valore di verità}.
È il caso di formule come $x < 0$, in cui la verità dipenderebbe dal «valore assunto da $x$», oltre che
dall'interpretazione dei simboli $<$ e $0$.

D'altra parte, se $p$ è un predicato unario, allora $p(x) \vee \neg p(x)$ sarà \textbf{sempre vera},
qualunque sia il valore di $x$ e qualunque sia il significato di $p$.

\dfn[valida]{Formula valida}{
    Una FBF si dice \textbf{valida} se è vera in tutti i modelli, cioè qualunque sia l'interpretazione dei simboli
    (l'insieme su cui variano le variabili e il significato di costanti, lettere funzionali e lettere predicative)
    e qualunque sia il valore assegnato alle sue variabili libere.
}

Un esempio di \textit{modello} è proprio l'aritmetica $(\mbN, +, \cdot, 0, s, \leq, =)$ della Sezione~«\nameref{sec:aritmetica}».
Le formule chiuse (Definizione~\ref{def:chiusa}) sono quelle che, fissato un modello, sono vere o false senza bisogno
di dare un valore alle variabili.

\info{Le regole del Capitolo 1 continuano a valere}{
    Le formule valide giocano il ruolo che nella logica proposizionale avevano le tautologie. In particolare
    (e sarà usato di continuo):
    \begin{itemize}
        \item se in una tautologia sostituiamo le variabili proposizionali con FBF del primo ordine, otteniamo una
              formula valida (come nella Proposizione~\ref{prop:sost-variabili}); ad esempio $\neg(\neg a) \leftrightarrow a$
              è valida per ogni FBF $a$;
        \item diremo che due FBF $a$ e $b$ sono \textbf{logicamente equivalenti} se $a \leftrightarrow b$ è valida; come nella
              Proposizione~\ref{prop:sostituzione}, sostituendo in una FBF una sottoformula con una equivalente
              (anche dentro lo scopo di un quantificatore) si ottiene una FBF equivalente;
        \item le equivalenze si possono concatenare come nel Corollario~\ref{cor:catene}.
    \end{itemize}
}

% ──────────────────────────────────────────────────────────────
\section{Negazione dei quantificatori}\label{sec:negazione-quant}

Come nego «Ogni gatto è bianco»? Dicendo che «\textbf{Esiste} un gatto che \textbf{non} è bianco».
Quindi negare un $\forall$ produce un $\exists$ (e viceversa), e la negazione «entra» nello scopo.

\begin{concetto}
\begin{Teorema}{Negazione dei quantificatori}{negazione-quantificatori}
    \begin{prerequisiti}
        \dip{Definizione}{def:esiste}
        \dip{Definizione}{def:valida}
        \dip{Proposizione}{prop:doppianeg}
        \dip{Proposizione}{prop:negdoppiaimp}
        \dip{Proposizione}{prop:sost-variabili}
        \dip{Proposizione}{prop:sostituzione}
        \dip{Corollario}{cor:catene}
        \dipsez{sec:valide}
    \end{prerequisiti}
    Per ogni FBF $a$ e ogni variabile $x$ le seguenti formule sono valide:
    \begin{enumerate}
        \item $(\forall x)(\neg a) \leftrightarrow \neg\big((\exists x)(a)\big)$;
        \item $\neg\big((\forall x)(a)\big) \leftrightarrow (\exists x)(\neg a)$;
        \item $\neg\big((\exists x)(\neg a)\big) \leftrightarrow (\forall x)(a)$.
    \end{enumerate}
\end{Teorema}

\begin{dimostrazione}
    \textbf{Punto 1.} Per la Definizione~\ref{def:esiste}, $(\exists x)(a)$ è per definizione $\neg\big((\forall x)(\neg a)\big)$.
    Quindi $\neg\big((\exists x)(a)\big)$ è la formula $\neg\neg\big((\forall x)(\neg a)\big)$, che per la doppia negazione
    (applicata alla FBF $(\forall x)(\neg a)$) è equivalente a $(\forall x)(\neg a)$.

    \textbf{Punto 2.} Applichiamo la Definizione~\ref{def:esiste} alla FBF $\neg a$:
    \[ (\exists x)(\neg a) \ \text{ è per definizione } \ \neg\big((\forall x)(\neg\neg a)\big). \]
    Per la doppia negazione $\neg\neg a \leftrightarrow a$; sostituendo dentro lo scopo del quantificatore
    (Proposizione~\ref{prop:sostituzione}) otteniamo $(\forall x)(\neg\neg a) \leftrightarrow (\forall x)(a)$ e quindi,
    per la negazione della doppia implicazione, $\neg\big((\forall x)(\neg\neg a)\big) \leftrightarrow \neg\big((\forall x)(a)\big)$.
    Dunque $(\exists x)(\neg a) \leftrightarrow \neg\big((\forall x)(a)\big)$, che è il punto 2 letto da destra a sinistra.

    \textbf{Punto 3.} Applicando la negazione della doppia implicazione al punto 2 otteniamo
    \[ \neg\neg\big((\forall x)(a)\big) \leftrightarrow \neg\big((\exists x)(\neg a)\big), \]
    e per la doppia negazione $\neg\neg\big((\forall x)(a)\big) \leftrightarrow (\forall x)(a)$.
    Concatenando le due equivalenze (Corollario~\ref{cor:catene}) si ottiene il punto 3.
\end{dimostrazione}
\end{concetto}

In pratica: \textbf{per negare un quantificatore si scambia $\forall$ con $\exists$ e si nega lo scopo}.

\begin{esempio}{Ogni gatto è bianco}
    Siano $G(x)$ = «$x$ è un gatto» e $B(x)$ = «$x$ è bianco». «Ogni gatto è bianco» si scrive
    $\forall x\,\big(G(x) \rightarrow B(x)\big)$. Negando:
    \[
    \begin{aligned}
        \neg\forall x\,\big(G(x) \rightarrow B(x)\big)
        &\leftrightarrow \exists x\,\neg\big(G(x) \rightarrow B(x)\big) && \text{(Teorema~\ref{teo:negazione-quantificatori}, punto 2)}\\
        &\leftrightarrow \exists x\,\big(G(x) \wedge \neg B(x)\big) && \text{(negazione dell'implicazione)}
    \end{aligned}
    \]
    cioè «esiste un gatto che non è bianco».
\end{esempio}

Il prossimo esempio fa parte della teoria, perché verrà usato più avanti.

\begin{esempio}{Negare $(\exists x)\big(f(x) \wedge \forall y(\neg g(x,y))\big)$}
    \begin{prerequisiti}
        \dip{Teorema}{teo:negazione-quantificatori}
        \dip{Definizione}{def:esiste}
        \dip{Proposizione}{prop:demorgan}
        \dip{Proposizione}{prop:impor}
        \dip{Proposizione}{prop:sostituzione}
        \dip{Corollario}{cor:catene}
    \end{prerequisiti}
    Qui $f$ e $g$ sono lettere predicative (unaria e binaria).
    \[
    \begin{aligned}
        \neg(\exists x)\big(f(x) \wedge \forall y(\neg g(x,y))\big)
        &\leftrightarrow (\forall x)\,\neg\big(f(x) \wedge \forall y(\neg g(x,y))\big) && \text{(Teorema~\ref{teo:negazione-quantificatori}, punto 1)}\\
        &\leftrightarrow (\forall x)\big(\neg f(x) \vee \neg\forall y(\neg g(x,y))\big) && \text{(De Morgan)}\\
        &\leftrightarrow (\forall x)\big(\neg f(x) \vee \exists y\,(g(x,y))\big) && \text{(definizione di $\exists$)}\\
        &\leftrightarrow (\forall x)\big(f(x) \rightarrow \exists y\,(g(x,y))\big) && \text{(relazione tra $\rightarrow$ e $\vee$)}.
    \end{aligned}
    \]
    A parole: «per ogni $x$, se $x$ ha la proprietà $f$, allora esiste $y$ tale che $g(x,y)$».
\end{esempio}

% ──────────────────────────────────────────────────────────────
\section{Sostituire termini alle variabili libere}\label{sec:sostituzione-termini}

\dfn[notazione-sostituzione]{Sostituzione di termini}{
    Se la formula $p$ ha come variabili libere $x_1, \dots, x_n$, possiamo scrivere $p(x_1, \dots, x_n)$.
    Se $t_1, \dots, t_n$ sono termini, scriviamo $p(t_1, \dots, t_n)$ per la formula $p$ con $t_1$ al posto di $x_1$,
    \dots, $t_n$ al posto di $x_n$.
}

Attenzione: scrivere $p(x_1, \dots, x_n)$ \textbf{non} vuol dire che queste siano \textbf{tutte} le variabili libere di $p$.
L'effetto della sostituzione è quello di mettere, al posto delle occorrenze libere delle variabili, dei termini «concreti».

\begin{esempio}{Sostituzione}
    Sia $p$ la formula
    \[ (x = y \wedge x \neq 0) \vee \exists z\,(x + z = y). \]
    Possiamo scrivere $p(x)$, come anche $p(x, y)$, ma \textbf{non} $p(x, z)$, perché ogni occorrenza di $z$ è vincolata.

    Prendiamo i termini $0$ e $1$ e sostituiamoli al posto di $x$ e $y$: $p(0, 1)$ diventa
    \[ (0 = 1 \wedge 0 \neq 0) \vee \exists z\,(0 + z = 1) \]
    ed è \textbf{vera} nell'aritmetica: la prima parte è falsa, ma la seconda è vera (basta prendere $z = 1$).
\end{esempio}

% ──────────────────────────────────────────────────────────────
\section{Convenzioni di scrittura}\label{sec:convenzioni}

\subsection{Quantificatori multipli}

Per brevità si usa raggruppare i quantificatori dello stesso tipo:
\[ \forall x_1 \Big(\forall x_2 \big(\cdots \forall x_n (p) \cdots\big)\Big) \quad \text{si scrive} \quad \forall x_1, x_2, \dots, x_n\,(p), \]
e lo stesso vale per $\exists$.

\subsection{Notazione infissa}

\dfn[infissa]{Notazione infissa}{
    Sia $p$ una lettera predicativa binaria e siano $t_1$, $t_2$ termini. Scriveremo
    \[ t_1 \, p \, t_2 \quad \text{invece di} \quad p(t_1, t_2), \]
    proprio come si fa per $\leq$ (scriviamo $x \leq y$ e non $\leq(x, y)$).
}

\subsection{Quantificatori ristretti}\label{sec:ristretti}

\dfn[ristretti]{Quantificatori ristretti}{
    Sia $p$ una lettera predicativa binaria, $x$ una variabile, $y$ un termine e $f$ una FBF. Scriveremo
    \[ (\forall x\, p\, y)(f) \ \defsse\ (\forall x)\big(p(x, y) \rightarrow f\big) \]
    \[ (\exists x\, p\, y)(f) \ \defsse\ (\exists x)\big(p(x, y) \wedge f\big) \]
    Al posto di $y$ può esserci un qualunque termine $t$.
}

A parole: $(\forall x\, p\, y)(f)$ vuol dire «per ogni $x$ che sta nella relazione $p$ con $y$, vale $f$» e
$(\exists x\, p\, y)(f)$ vuol dire «esiste un $x$ che sta nella relazione $p$ con $y$ e per cui vale $f$».

\begin{esempio}{Quantificatori ristretti nell'aritmetica}
    \begin{itemize}
        \item $(\forall x < s(0))(f)$ significa $\forall x\,\big(x < s(0) \rightarrow f\big)$: «per ogni $x$ minore di $1$ vale $f$».
        \item La formula $(\forall x < 0)(x \neq x)$ è \textbf{vera}. Perché? Significa $\forall x\,(x < 0 \rightarrow x \neq x)$
              e in $\mbN$ nessun $x$ è minore di $0$: la premessa $x < 0$ è sempre falsa, quindi l'implicazione è sempre vera
              (Proposizione~\ref{prop:tavole-connettivi}), per ogni $x$. Si dice che è vera «a vuoto»: non ci sono
              controesempi perché non c'è nessun $x$ da controllare.
        \item La formula $(\exists x < 0)(f)$, interpretata nell'aritmetica su $\mbN$, è invece \textbf{falsa} qualunque sia $f$:
              significa $\exists x\,(x < 0 \wedge f)$ e non esiste alcun $x$ con $x < 0$.
    \end{itemize}
\end{esempio}

\warning{{Con $\forall$ si usa $\rightarrow$, con $\exists$ si usa $\wedge$}}{
    È un errore comune scambiarli.
    \begin{itemize}
        \item $(\forall x)\big(p(x,y) \wedge f\big)$ direbbe che \textbf{ogni} $x$ sta nella relazione $p$ con $y$ e soddisfa $f$:
              molto più forte di quello che vogliamo.
        \item $(\exists x)\big(p(x,y) \rightarrow f\big)$ sarebbe vera appena esiste un $x$ che \textbf{non} sta nella
              relazione $p$ con $y$ (premessa falsa, implicazione vera): molto più debole di quello che vogliamo.
    \end{itemize}
    Ad esempio «alcuni psichiatri sono baristi» si scrive $\exists x\,\big(P(x) \wedge B(x)\big)$, non $\exists x\,\big(P(x) \rightarrow B(x)\big)$.
}

\begin{concetto}
\begin{Proposizione}{Negazione dei quantificatori ristretti}{neg-ristretti}
    \begin{prerequisiti}
        \dip{Definizione}{def:ristretti}
        \dip{Teorema}{teo:negazione-quantificatori}
        \dip{Proposizione}{prop:negimp}
        \dip{Proposizione}{prop:demorgan}
        \dip{Proposizione}{prop:impor}
        \dip{Proposizione}{prop:sostituzione}
        \dip{Corollario}{cor:catene}
    \end{prerequisiti}
    Con le notazioni della Definizione~\ref{def:ristretti}, le seguenti formule sono valide:
    \begin{enumerate}
        \item $\neg(\forall x\, p\, y)(f) \leftrightarrow (\exists x\, p\, y)(\neg f)$;
        \item $\neg(\exists x\, p\, y)(f) \leftrightarrow (\forall x\, p\, y)(\neg f)$.
    \end{enumerate}
\end{Proposizione}

\begin{dimostrazione}
    \textbf{Punto 1.}
    \[
    \begin{aligned}
        \neg(\forall x\, p\, y)(f) &\leftrightarrow \neg(\forall x)\big(p(x,y) \rightarrow f\big) && \text{(definizione)}\\
        &\leftrightarrow (\exists x)\,\neg\big(p(x,y) \rightarrow f\big) && \text{(Teorema~\ref{teo:negazione-quantificatori}, punto 2)}\\
        &\leftrightarrow (\exists x)\big(p(x,y) \wedge \neg f\big) && \text{(negazione dell'implicazione)}\\
        &\leftrightarrow (\exists x\, p\, y)(\neg f) && \text{(definizione)}.
    \end{aligned}
    \]
    \textbf{Punto 2.}
    \[
    \begin{aligned}
        \neg(\exists x\, p\, y)(f) &\leftrightarrow \neg(\exists x)\big(p(x,y) \wedge f\big) && \text{(definizione)}\\
        &\leftrightarrow (\forall x)\,\neg\big(p(x,y) \wedge f\big) && \text{(Teorema~\ref{teo:negazione-quantificatori}, punto 1)}\\
        &\leftrightarrow (\forall x)\big(\neg p(x,y) \vee \neg f\big) && \text{(De Morgan)}\\
        &\leftrightarrow (\forall x)\big(p(x,y) \rightarrow \neg f\big) && \text{(relazione tra $\rightarrow$ e $\vee$)}\\
        &\leftrightarrow (\forall x\, p\, y)(\neg f) && \text{(definizione)}.
    \end{aligned}
    \]
\end{dimostrazione}
\end{concetto}

In pratica: per negare un quantificatore ristretto si scambia $\forall$ con $\exists$, \textbf{si lascia invariata la restrizione}
e si nega lo scopo. Ad esempio $\neg(\forall x \leq 12)(f) \leftrightarrow (\exists x \leq 12)(\neg f)$.
